%!TEX root = ../nauchka.tex

\section{Реализация}

\begin{frame}
\frametitle{Архитектура AlphaBetaCROWN + auto\_LiRPA}

\begin{block}{auto\_LiRPA}
Библиотека для автоматического вычисления границ (Linear Relaxation based Perturbation Analysis)
\end{block}

\vspace{0.3cm}

\textbf{Ключевые компоненты:}
\begin{itemize}
    \item \texttt{BoundedModule} --- обертка над PyTorch моделью
    \item \texttt{Bound} операторы для каждого типа слоя
    \item Метод \texttt{bound\_backward} --- обратное распространение границ
\end{itemize}

\vspace{0.3cm}

\begin{exampleblock}{Идея реализации}
Создать новые классы операторов \texttt{BoundLinearTP\_Col} и \texttt{BoundLinearTP\_Row}, которые переопределяют \texttt{bound\_backward} для распределенных вычислений
\end{exampleblock}

\end{frame}

\begin{frame}[fragile]
\frametitle{Что было сделано: Структура кода}

\textbf{Созданные файлы:}

\begin{enumerate}
    \item \texttt{auto\_LiRPA/operators/tensor\_parallel.py}
    \begin{itemize}
        \item \texttt{BoundLinearTP\_Col} --- Column Parallel оператор
        \item \texttt{BoundLinearTP\_Row} --- Row Parallel оператор
    \end{itemize}

    \item \texttt{examples/simple/tp\_model.py}
    \begin{itemize}
        \item \texttt{ColumnParallelLinear} --- PyTorch слой с разделенными весами
        \item \texttt{RowParallelLinear} --- PyTorch слой с разделенными весами
    \end{itemize}

    \item \texttt{examples/simple/test\_tp\_verification.py}
    \begin{itemize}
        \item Тестовый скрипт для демонстрации распределенной верификации
    \end{itemize}
\end{enumerate}

\vspace{0.3cm}

\textbf{Всего:} 3 новых модуля, $\sim$400 строк кода

\end{frame}

\begin{frame}[fragile]
\frametitle{Реализация: BoundLinearTP\_Col}

\begin{block}{Псевдокод Column Parallel слоя}
\begin{verbatim}
def bound_backward(self, last_lA, last_uA, ...):
    # Вызвать базовый метод для локальных вычислений
    result = super().bound_backward(...)

    # Извлечь A-матрицы
    lA_x, uA_x = result[0][0]

    # AllReduce для объединения результатов
    dist.all_reduce(lA_x, op=ReduceOp.SUM)
    dist.all_reduce(uA_x, op=ReduceOp.SUM)

    return result
\end{verbatim}
\end{block}

\textbf{Ключевая идея:} Использовать \texttt{torch.distributed.all\_reduce} для суммирования частичных произведений

\end{frame}

\begin{frame}[fragile]
\frametitle{Реализация: BoundLinearTP\_Row}

\begin{block}{Псевдокод Row Parallel слоя}
\begin{verbatim}
def bound_backward(self, last_lA, last_uA, ...):
    # Вызвать базовый метод
    result = super().bound_backward(...)

    # Результат уже корректно разделен!
    # Коммуникация НЕ нужна

    return result
\end{verbatim}
\end{block}

\vspace{0.3cm}

\textbf{Особенность:} Row Parallel слои автоматически производят разделенный результат без дополнительной коммуникации

\end{frame}

\begin{frame}
\frametitle{Статус реализации}

\begin{block}{Что работает}
\begin{itemize}
    \item[\checkmark] Код компилируется и интегрируется с auto\_LiRPA
    \item[\checkmark] Базовый пример верификации работает на 2 GPU
    \item[\checkmark] Классы TP операторов реализованы
    \item[\checkmark] Тестовая модель с TP слоями создана
\end{itemize}
\end{block}

\vspace{0.3cm}

\begin{alertblock}{Что требует доработки}
\begin{itemize}
    \item Полное тестирование на 2+ GPU
    \item Тестирование на больших моделях
    \item Интеграция с полным циклом AlphaBetaCROWN
    \item Оптимизация коммуникаций (overlapping)
\end{itemize}
\end{alertblock}

\end{frame}
